\documentclass[11pt,a4paper]{article}
\usepackage[T1]{fontenc}
\usepackage{lmodern}
\usepackage[margin=22mm]{geometry}
\usepackage{amsmath,amssymb,amsthm,microtype}
\usepackage[hidelinks]{hyperref}
\setlength{\parindent}{0pt}
\setlength{\parskip}{5pt}
\setlength{\emergencystretch}{2em}
\newtheorem{lemma}{Lemma}
\title{The Nine-Dimensional Snake Problem:\\A Short Explanation of the Current Argument}
\author{Iori Hironaka \quad Tommy Ma \quad Steven Zhan}
\date{Research state after round 239}
\begin{document}
\maketitle

\section{Question and evidence}
Let $a(n)$ be the largest edge length of an induced path in the binary
hypercube $Q_n$. Its vertices are distinct, and two are adjacent exactly
when their path indices are consecutive. A length-$m$ snake has $m+1$ vertices.

The maintained interval is
\[
\boxed{191\le a(9)\le231.}
\]
The lower bound comes from a verified construction. The contributor's
completed 231 theorem and saved final assembly are pinned to revision
\texttt{478869ee}. Its standard-axiom audit passed, and no certificate or
share-floor premise remains in the theorem. A fresh independent full
rebuild remains pending~\cite{bound231archive}.

The known values $1,2,4,7,13,26,50,98$ do not determine the next term.
Numerical continuations suggest 194, 195 or 196, without proving a recurrence.
To establish 196 exactly requires both a 196-edge witness and exclusion
of every 197-edge snake; every longer snake contains such a prefix.
The counting argument supplies the upper bound. The construction arguments
seek any improvement above 191, beginning with a 192-edge snake. Their
restricted exclusions do not establish the global optimum.

The current reconstruction argument recovers a binary selection from
any 508 labelled neighbour counts. It also permits up to eight missing
counts when no pair of missing positions is at distance eight. These
are written cube deductions with checked arithmetic; the fully assembled
Lean cube theorem currently covers three missing counts.

\section{The square budget and the completed upper bound}
For a nontrivial snake put $v=m+1$, and let $s_j$ count coordinate squares
containing exactly $j$ selected vertices. Write $T=s_1$.
There are 4608 squares, each selected vertex lies in 36, and no square
is fully selected. Three-selected-corner squares correspond exactly to
internal path vertices, so $s_3=v-2$. Consequently
\[
s_0+s_1+s_2+s_3=4608,\qquad s_1+2s_2+3s_3=36v,
\]
and subtraction gives the exact identity
\[
\boxed{T+2s_0+35(m+1)=9214.}
\]

If $\alpha_i$ counts singleton squares at path vertex $i$, local geometry
gives positive counts, endpoint counts at least four, internal neighbours
of a count one at least four, and no three consecutive internal twos.
These restrictions imply sums at least 14 and 16 on six and seven
consecutive internal vertices. For example, every triple sums to at least
six, with equality only at $(1,4,1)$; the neighbouring count restriction
handles that exceptional triple when combining blocks.
The archived count-only potential and omitted boundary counts yield
$T\ge4m-13$, hence $39m+2s_0\le9192$. This historical route gives 235;
its potential retains a separate independent-rebuild requirement.

Empty squares retain information discarded by that count-only estimate.
For a cube vertex $q$, let $e(q)$ count its incident empty squares and
$d(q)$ its selected neighbours. Each selected vertex $x$ receives
\[
\sigma_{12}(x)=\sum_{q\sim x}
 \left\lfloor\frac{12e(q)}{d(q)}\right\rfloor.
\]
All denominators appearing here are positive. Let $L$ sum the remainders
$12e(q)\bmod d(q)$ for $d(q)>0$, and $12e(q)$ for $d(q)=0$.
Division with remainder and the four corners of every empty square give
\[
\sum_{x\in S}\sigma_{12}(x)+L=48s_0,\qquad
G+840(m+1)+L=221136,
\]
where $S$ is the selected vertex set and $G=24T+\sum_{x\in S}\sigma_{12}(x)$.

In the earlier completed 233 comparison, its total window weight $W$
satisfies, for $m\ge7$,
\[
520(m-7)\le W+1023,\qquad W+4200m\le1099800.
\]
The checked local certificates, share floors and finite potential therefore give
\[
4720m\le1104463<4720\cdot234=1104480,
\]
so $m\le233$. Shorter paths already satisfy this bound. The saved assembly
reuses imported objects and cloud-built certificates; their fresh independent
rebuild remains distinct from the successful saved audit~\cite{bound233archive}.

The later contributor theorem at revision \texttt{478869ee} supplies the
current upper bound 231. Its final argument includes the required local
refutations and share floors. The saved final assembly and standard-axiom
audit passed; a fresh independent whole-proof rebuild remains pending
~\cite{bound231archive}. The displayed 233 calculation is retained as the
earlier argument and is not used to derive 231.

The alternative retained comparison uses $m-8$ windows. Under its own
local-weight and state-walk hypotheses, boundary counting and telescoping give
\[
(4200+\mu)m+5L\le1098600+8\mu+R,
\]
where every state edge has weight plus potential drop at least $\mu$,
and the potential has range at most $R$. Its conditional 234 result,
and a separate relaxed 233 comparison, still require their stated local
certificates and geometric and graph interpretations. Neither improves the
maintained interval. Their unresolved premises do not belong to the completed
contributor theorem; detailed profiles remain in the companion manuscript.

\section{Joining lower-dimensional paths}
Cutting at the $k$ edges in one coordinate leaves $k+1$ path components
in the two faces, with total internal length $m-k$.
The saved representatives have decompositions
\[
26=2+13+9+2,\quad 50=18+24+6+2,\quad
98=8+25+10+25+26+4,
\]
where each final term counts crossings. These describe particular witnesses.

Paths of lengths $p,q$ in opposite faces join into a snake of length
$p+q+1$ if their projections meet only at their designated joining endpoint.
Cross-face adjacency requires identical projections, which proves the claim.
For a two-crossing construction, take disjoint, nonadjacent outer paths
$P,R$ in one face and a middle path $Q$ in the other. Each projected
intersection must consist of one endpoint of the corresponding outer path
and a different endpoint of $Q$. The resulting length is
$|E(P)|+|E(Q)|+|E(R)|+2$, explaining the checked $18+24+6+2=50$ example.

There is a restriction on modifying two positive-length parents $p,q$.
Delete $D$ original vertices and add $A$ new ones. Let $b$ count deleted
original endpoints, $e$ original edges deleted at both ends, and $H$ final
within-face edges touching a new vertex. If the result has length $m$
and $k$ crossings, vertex and edge counting give
\[
m+k+b+e+H=p+q+2+2A,\qquad m+k\le2a(n)+2+2A.
\]
Indeed, deletion removes $2D-b-e$ original edges, and the final vertex
count is $p+q+2-D+A=m+1$.
Without additions, a 196-edge lift from parents of length at most 98
uses lengths 98 and 97 without deletion, or two lengths 98 with one
deletion; it permits at most two crossings. These are necessary conditions.

The saved 98-edge $Q_8$ representative crosses every coordinate at least
four times. Its only four-crossing split has two 25-edge components in
one face, using 52 vertices, whereas a single $Q_7$ snake has at most 51.
The counting restriction excludes all other splits from a deletion-only
two-parent lift. This representative therefore already requires a more
general construction than deletion from two $Q_7$ parents.

\section{Changing the embedding of an entire snake}
Partition coordinates and write vertices as $(z,u)$. Assign an isometry
$g_z$ of the free-coordinate cube to each retained slice, and map
$(z,u)$ to $(z,g_z(u))$. Within-slice distances are preserved, and
nonadjacent slices cannot produce cube edges.
For adjacent slices with selected free-word sets $B_z,B_w$, put
$h=g_z^{-1}g_w$. All remaining edges and nonedges are preserved exactly when
\[
h(u)=u\quad(u\in B_z\cap B_w),\qquad
h(B_w\setminus B_z)\cap B_z=\varnothing.
\]
The conditions retain old crossings and forbid new projected coincidences.

This family produces a different cube class from the saved 50-edge parent.
For the saved 98-edge parent, the complete native census over all 254
nontrivial coordinate partitions produces no new class. Lean checks the
generic criterion and coverage of the literal finite tables. Their complete
geometric interpretation and the surviving leaves retain native status.

With two free coordinates, a written classification uses a uniform parity
offset followed by component switches; its complete Lean assembly is pending.
With three free coordinates, an injective partial map preserving parity,
adjacency and nonadjacency extends to an isometry: it permutes antipodal-pair
labels consistently. This extension fails in four coordinates, where equal
indexed parity patterns can have different endpoint distances. Arbitrary
coordinate histories therefore allow changes beyond slice isometries.

\section{Coordinate histories and finite capacity}
Retain the source base word $z_i$ and choose arbitrary free words $u_i$.
For Hamming distance $d$, the exact path conditions are
\[
d(z_i,z_j)+d(u_i,u_j)
\begin{cases}=1,&j=i+1,\\ \ge2,&j\ge i+2.\end{cases}
\]
They exclude both repeats and chords. Consecutive retained-coordinate steps
force a constant free word, so merge their maximal index intervals while
retaining every pair constraint. Each interval change flips one free bit;
interval word parities alternate. Translation, coordinate permutation and
lexicographic representatives remove symmetry without losing an orbit.

For this indexed 98-edge source, all 70 four-free-coordinate and all 56
five-free-coordinate models exhaust after source-isometric assignments are
excluded. This native result concerns the specified source and index matching.
For six free coordinates, retained pair $(1,4)$ also exhausts; 27 choices
remain unresolved. Four have saved unfinished searches, at $(0,4)$,
$(3,4)$, $(0,5)$ and $(0,6)$; the other 23 have only limited initial runs.
The latest $(0,6)$ continuation retains eight unfinished regions and no new parent.

With two retained coordinates, the base is a square. Distinct intervals of
the same free-word parity cannot have consecutive original indices.
If they share a word, every cross-interval base pair must be opposite.
A square corner has a unique opposite, so path injectivity forces both
intervals to be singletons at opposite corners. No three can share one word.

For one parity, let $M$ count intervals visiting multiple base vertices and
$N_z$ singleton intervals at corner $z$. Their distinct-word capacity gives
\[
\boxed{M+\max(N_{00},N_{11})+\max(N_{01},N_{10})\le2^{r-1}}
\]
for $r\ge1$ free coordinates: choose all multi-visit intervals and the larger
singleton family from each opposite pair. Every two chosen words differ.
For retained pair $(0,5)$ the even counts are $M=12$ and
$(N_{00},N_{01},N_{10},N_{11})=(5,4,10,10)$.
The chosen 32 intervals use every even word of $Q_6$; each of the nine
outside intervals must share with a chosen antipodal partner.

For a distinct-word family and a subset $J$, let $D_j$ be current domains
and $U=\bigcup_{j\in J}D_j$. Necessarily $|J|\le|U|$.
At equality, injectivity uses every word of $U$ within $J$, so other
members of that family must avoid $U$. Such removal requires the stated
distinctness and cannot automatically be applied across partner families.

The outside singleton family also has distinct words: its corners are equal
or adjacent. If $C$ is the chosen family and $O$ the outside family of one
parity in $Q_6$, at most $s=32-|C|$ outside values are absent from $C$.
For $(0,6)$, $s=2$ in even parity and $s=6$ in odd parity. Thus at least
seven of nine even and seven of thirteen odd outside intervals must share.
An outside interval with no domain intersection with any permitted partner
must consume one spare word. Once $s$ such intervals are forced, every
other outside interval must retain support in its partners' domains.

More generally, let $N(J)$ contain the chosen partners of outside intervals
$J$ whose domains intersect. Shared members inject into distinct partners,
so $|J|\le|N(J)|+s$. At equality, $J$ uses every value assigned to $N(J)$
and exhausts the spare allowance. Other outside intervals must therefore
share with chosen partners outside $N(J)$.

Lean checks the generic counting and assignment-preservation arguments.
For the literal $(0,6)$ source, checked index witnesses connect the actual
path to chosen/outside distinctness and complete partner lists. The normalized
initial word, explicit six-bit range and group boundaries also yield the
two 32-word parity pools. Current domain fitting remains an explicit premise;
native relation tables, matching discovery and complete traversal remain
separate from these checked path and filtering implications.

Consecutive triples supply another restriction. Two free steps leave the
outer words at distance zero or two, and each retained pair needs a compatible
third value in its current domain. For example, words 5 and 6 have only
common neighbours 4 and 7; removing both excludes that pair.
Lean checks preservation of every fitting assignment and validity under
smaller domains. The native continuation retains projected pair relations
within a job and reloads original relations for each new job.
No controlled comparison establishes a speed improvement.

Continuation itself relies on coverage: splitting $D$ into $\{a\}$ and
$D\setminus\{a\}$ partitions assignments. Saving unfinished descendants and
unexamined ancestor values preserves the unexplored region. Lean checks this
generic implication; saved-state audits do not certify the entire native search.

\section{The completed fixed-parent exclusion}
Let $P$ be the saved 98-edge $Q_8$ parent and $t$ either endpoint, 0 or 78.
A 93-edge induced path beginning at $t$ and avoiding $V(P)\setminus\{t\}$
would give $98+93+1=192$ edges after one crossing. Each permitted graph
has 158 vertices and 337 edges. Earlier degree-and-cycle searches returned
\texttt{UNKNOWN}; the prefix certificates below now exclude both targets.

A 93-edge partner has 47 vertices of each parity and an odd free endpoint.
The model gives degree one to the two selected endpoints and degree two
to other selected vertices, with $y_{ab}=x_ax_b$ for Boolean selections.
An induced path omits a vertex of every cycle: its last cycle vertex
cannot have two distinct earlier cycle neighbours, since both would be
its immediate predecessor. Lean checks the parity and cycle implications.

For every allowed subset $U$ and $v\in U$, selected edges also satisfy
\[
\sum_{ab\in E(U)}y_{ab}\le\sum_{u\in U}x_u-x_v.
\]
Here $E(U)$ contains the allowed edges within $U$. Inducedness assigns
each selected edge to a distinct selected path index, avoiding a selected
marked index. This proves the inequality; the unselected case is weaker.
If $X,Y,B,E$ count selected vertices, internal edges, boundary edges and
selected endpoints in $U$, degree accounting gives $2X-E=2Y+B$.
Thus the same inequality is $B+E\ge2x_v$. These deductions pass Lean.

For Boolean products, $\delta_{ab}=x_a+x_b-2y_{ab}$ records endpoint
disagreement. A closed cycle $C$ changes value an even number of times.
An odd marked edge set $F$ must differ from that pattern somewhere, giving
\[
\sum_{e\in F}\delta_e-\sum_{e\in E(C)\setminus F}\delta_e\le|F|-1.
\]
Lean checks the arbitrary closed-walk argument. These inequalities are
classical~\cite{barahona1986}. The preceding finite comparison gave exact
cardinality bounds 95.955197 and 95.763979, excluding 95-edge partners.

\paragraph{Why the continuous model alone remains insufficient.}
The exact fractional feasible cardinalities are
\[
q_0=\frac{2507137819890885477}{26128236866583647},\qquad
q_{78}=\frac{388107824092577}{4052758842992}.
\]
Both exceed 94 and satisfy all three complete inequality families.
For occupancy cycles, put $w_u=1-x_u$. For each deleted edge $ab$,
a potential satisfies $p_a=0$, $|p_u-p_v|\le w_u+w_v$ on other edges,
and $p_b+w_a+w_b\ge2$. Summing around any cycle gives
$\sum_{u\in C}w_u\ge1$. For subset inequalities, 158 feasible flows
per root send at least $2x_v$ from $v$ to an added sink. Conservation
bounds the supply by every separating cut, whose capacity is $B+E$.

For parity cycles, make two copies of each vertex. Edges preserving
the copy label cost $\delta_{ab}$; switching edges cost $1-\delta_{ab}$.
For each vertex, a saved potential increases by at least one between
its copies and satisfies $p(b)-p(a)\le c_{ab}$ on every directed arc.
A cycle switching on an odd marked set ends at the opposite copy.
Telescoping gives cost at least one, the required parity inequality.
Independent exact arithmetic audits verify these rational potentials and flows.

Mix either point with a one-edge path, using weight $92/(q_t-2)$ on
the larger point. The mixture has weight 94 and preserves all three
linear families. Thus this continuous relaxation still admits the target;
these fractional points supply no path~\cite{research220,research221}.

\paragraph{Complete division by actual prefixes.}
For a hypothetical path $q_0,\ldots,q_{93}$, each proper prefix with
$k<93$ requires $x_{q_i}=1$ for $0\le i\le k$ and $e_{q_i}=0$ for
$0<i\le k$. Injectivity makes every nonroot prefix vertex different
from both full endpoints, including the current last prefix vertex.
Its next vertex is allowed, adjacent to the last vertex, unused, and
nonadjacent to every earlier prefix vertex. Enumerating all such vertices
covers every possible continuation. Lean derives these necessities.

Each branch adds its prefix equalities to the model. Each exact
certificate combines inequality rows with nonnegative weights and equality
rows with signed weights, then bounds the remaining coefficients using
$0\le z_j\le1$; numerical optimality is unnecessary. Cardinality bounds
below 94 exclude branches; every other branch is expanded completely.
The audited counts, with depth measured in prefix edges, are
\[
\begin{array}{c|rrrrrr|r}
\text{Depth}&1&2&3&4&5&6&\text{Total}\\\hline
\text{Nodes}&14&33&21&9&1&2&80\\
\text{Excluded}&3&23&17&8&0&2&53\\
\text{Expanded}&11&10&4&1&1&0&27
\end{array}
\]
There are no pending leaves. Root 0 has 31 excluded leaves and root 78
has 22; their largest exact bounds are 93.963070 and 93.984077.
The unique depth-four survivor was $(78,206,198,196,197)$, followed
necessarily by 213 and then by either 215 or 221. The final bounds are
93.979475 and 93.414818. The initial low-memory interruption preserved
45 completed new queries and 20 pending prefixes; continuation completed
the cover. The failed batch remains recorded separately.

Every actual 93-edge partner must follow this complete tree to a leaf
of depth at most six. Its 94 selected vertices contradict that leaf's
bound below 94; all used prefixes are proper because $6<93$.
Every longer partner has a compatible 93-edge prefix. Therefore
\[
\boxed{\text{partner length}\le92,\qquad
\text{one-crossing joining length}\le98+92+1=191.}
\]
An inverse cube isometry preserves length, inducedness, the joining
endpoint and parent avoidance, so the ceiling holds for the full
cube-isometry class of this parent. This pullback is also checked in Lean.
No compatible 92-edge witness or attainment of 191 by this construction
is established. Different parent classes and other face arrangements
are outside the exclusion.

Four generic results in two Lean modules support the prefix and isometry
arguments. Complete finite coverage and numerical certificates have
independent exact audits, with concrete model necessity reviewed separately.
The numerical tree and coefficient arrays have no end-to-end Lean assembly.
All evidence is retained in milestone 222~\cite{research222}.

\section{Joint construction across coordinate faces}
The initial four-crossing route places three paths in one $Q_8$ face
and two in the other; the later replacement search also permits other
crossing counts. It follows the saved lower-dimensional patterns
\[
18+24+6+2=50,\qquad 8+25+10+25+26+4=98.
\]
The second pattern places 52 vertices in one face as two paths, more
than a single longest $Q_7$ path can occupy. This is why both faces and
their joining endpoints should be optimized together.

Five oriented paths in alternating faces join to an induced path when
consecutive projections meet only at their joining endpoint,
nonconsecutive same-face components are disjoint and nonadjacent, and
nonconsecutive opposite-face projections are disjoint. Indeed, a crossing
edge exists exactly when two projections coincide; these conditions
leave precisely the intended path edges. The resulting length is the
sum of the five lengths plus four. This assembly constructor has now
passed Lean under those premises; existence of large compatible
components is still the main difficulty.

The chosen 191-edge seed has component lengths $46,83,32,13,13$.
Both face selections and all four crossing locations are variable in a
search for 192 edges. Endpoint 0 is retained and the other endpoint is
free. Parity forces the new endpoint to be even, and four crossings
place it in the starting face. Keeping both old endpoints would be
incompatible with the target length.

The native query allowing at most four seed omissions, with at least
one in each face, is infeasible. Thus a successful change in this
restricted class needs at least five omissions and six additions.
This exclusion has no Lean contradiction certificate. The eight- and
twelve-omission queries are unresolved at their recorded time limits;
neither provides an impossibility result. No longer snake has been
obtained. The checked joining theorem and the computational records
have separate scopes~\cite{research223}.

An endpoint-directed variant retains the old endpoint 88 internally
and seeks endpoint 80. Since both 88 and the old crossing vertex 208
are adjacent to 80, vertex 208 must be removed. Otherwise the proposed
endpoint has two selected neighbours. The general endpoint exclusion
argument has passed Lean. The corresponding twelve-omission search
remains unresolved after 75 seconds. This gives a concrete way to
force a crossing change without fixing the rest of either face.

\subsection{Why simultaneous changes must interact}
Suppose a family is closed under individual length-preserving block
replacements of at most $K$ edges, and none of its paths can be extended.
A simultaneous change of two such blocks with disjoint interiors can
enable extension only if each new interior contains a vertex equal or
adjacent to an original interior vertex of the other block. Otherwise the noninteracting replacement can be
performed first, followed by the other; closure would keep the final
path in the blocked family. The actual intermediate-path construction
and this conditional implication have now passed Lean.

Appending endpoint 24 to the chosen seed requires changing positions
7 and 143, whose original vertices are 28 and 56. All pairs of blocks
of at most 33 edges containing those positions are covered by 224
maximal window pairs. Geometry excludes 171 and native solver queries
exclude the other 53, even without imposing the interaction condition.
Thus this particular route is computationally closed without relying
on the earlier closure census. Its solver contradictions are not checked
in Lean.

The initial search with 44-edge blocks considered 301 maximal contexts:
107 were geometrically excluded, 102 were solver-excluded and 92 were unresolved. Connectivity
flows did not resolve 16 selected retries. These results concern the
fixed seed, fixed outside indices, endpoint 24 and four crossings.
They give no global optimality conclusion; larger or differently
arranged changes remain possible~\cite{research224}.

\subsection{The current replacement frontier}
Two local replacement paths construct a full snake when each is induced,
each new interior is disjoint and nonadjacent to retained vertices
outside its own closed interval, and the
two new interiors are disjoint and nonadjacent. Every pair of vertices
then lies in one local path, is retained, or is separated by one of these
conditions. This proves distinctness and exactly consecutive adjacency.
The direct constructor has passed Lean.

Enlarging a replacement interval keeps the same complete candidate:
the extra positions retain their original vertices. An excluded
containing context therefore excludes the smaller one with the same
target. The outside-agreement implication has passed Lean; numerical
exclusions still rely on the saved native calculations.

The four-crossing construction is now excluded through two replacement
blocks of 46 edges. All 315 maximal contexts are covered: 106 by initial
geometry, 185 by fixing the upper pieces, and 24 by direct search. In
the reduced group, 89 contexts have no compatible upper configuration.
The other 96 reduce to 496 path-pair tasks: one is contained in an earlier
excluded task and 495 are newly excluded. Eight initial timeouts required
longer runs. These remain native computational conclusions. The
four-crossing reduction uses the preceding native width-45 cover as the
premise of its boundary-change condition.

The exact reason for the projection restriction is simple. Across the
two layers, vertices are adjacent precisely when all other coordinates
agree. A lower path must avoid projections of its own upper pieces
except at its actual joining endpoints, and must avoid all projections
of the other replacement's upper pieces. A fixed endpoint that does not
join the upper layer has no exemption. For a single crossing,
the intended endpoint equality and injectivity of both paths make this
restriction equivalent to the exact joining condition. Lean checked
that equivalence and the positive-length fixed-endpoint consequence.

There is also a general crossing-allocation argument. Let $A,B$ be the
numbers of upper anchors of the replacements, each between zero and two,
and let $a,b$ be their crossing counts. Endpoint faces give
$a\equiv A\pmod2$ and $b\equiv B\pmod2$. If the upper face is too small
to hold a replacement with two upper anchors, zero crossings is excluded.
Then $a\ge A$ and $b\ge B$: a target of two requires a positive even
count, and a target of one requires an odd count. If the total budget
is $A+B$, both inequalities must be equalities. Lean checked this
arithmetic implication for arbitrary natural counts. The finite input
conditions still require their separate domain calculation.

To investigate whether fixing four crossings was obstructing construction,
the same 315 contexts were examined with more freedom. The six-crossing
query excluded 200 of the 209 nonstatic cases and left nine unresolved.
The next model left the total crossing count free. It kept the same
seed, endpoint 24, two block lengths and retained outside vertices, and
used no boundary restriction from the earlier four-crossing family.

The earlier results still provide valid count exclusions in each matching
context: total four is impossible, and total six is impossible in the
200 cases with completed six-crossing results. Adding these exclusions
to the larger model removes only cases already covered by native evidence.
At the end of round 229, the search with unrestricted crossing count
excluded 304 of the 315 maximal contexts. All nine earlier six-crossing
unknowns were excluded by broader queries that had imposed no six-count
restriction on those nine. The exact-six family was therefore covered
without circular reasoning~\cite{research229}.

In the eleven residual contexts, two total crossings force replacement
allocations $(0,0)$ in four cases and $(1,0)$ in seven. Enumerating the
short upper pieces gives fifteen pairs of paths to find in $Q_8$.
All fifteen native queries are infeasible. This excludes two crossings
throughout the selected family.

\section{Why repeated face visits require space}
Split a snake at every crossing of one chosen coordinate. Each remaining
piece is a consecutive visit to one of the two faces. An internal visit
has a crossing before and after it. It needs at least three vertices:
with one vertex, the two flanking vertices would coincide; with two,
the flanking vertices would be adjacent at nonconsecutive positions.
Both violate the snake condition.

Suppose there are $c>0$ crossings, $u$ vertices in the upper face and
$r$ endpoints there. If the path makes $k$ upper visits, exactly $r$
are endpoint visits. Counting their incident crossings and vertices gives
\[
c=2k-r,\qquad u\ge3(k-r)+r.
\]
Eliminating $k$, and then repeating the argument for the lower face, gives
\[
\boxed{3c\le2u+r,\qquad 3c\le2l+2-r.}
\]
Here $l$ counts lower vertices. The inequalities are immediate when
$c=0$. For a path of $m$ edges, adding them gives $3c\le m+2$.
A 46-edge replacement therefore crosses any chosen coordinate at most
16 times. Lean checked the local obstruction and the finite counting
implications. Extracting the visits from an arbitrary whole path remains
a written proof.

Every full candidate in the eleven residual cases has endpoints in the
same face and retains a crossing. Its total crossing count is therefore
positive and even. Since two, four and six are excluded, it must have at
least eight. The last search adds this requirement and the face-visit
inequalities while allowing every remaining crossing count.
All eleven queries are infeasible. Together with the earlier results,
this excludes all 315 maximal contexts and, by padding, every shorter
replacement in the same two-block family. This remains a native
computational exclusion with explicit seed, endpoint, anchor and
length-preservation restrictions~\cite{research230}.

\section{Reorder one colour class freely}
Cube vertices have even or odd parity according to their number of one
bits. Neighbours always have different parity. Choose a set $U$ of one
parity. Each opposite-parity vertex with exactly two neighbours in $U$
can act as a connection between those two vertices.

Make a graph of these connections on $U$. A path visiting every vertex
of this graph exactly once, a Hamiltonian path, expands into a snake by
inserting the corresponding connection vertex between each pair.
There are no chords: each inserted vertex has precisely its two intended
neighbours in $U$, and same-parity vertices are never adjacent.
Conversely, every snake with both endpoints in $U$ and this exact parity
class gives such an ordering. The resulting length is $2|U|-2$.
Lean checked the general interleaving construction and converse adjacency
and distinctness statements. The finite cube-data extraction remains separate.

For 192 edges, $|U|=97$. The 131 saved 191-edge snakes have 145 distinct
96-vertex parity classes. Adding every possible new same-parity vertex
gives 23,127 distinct candidate sets. Simple degree obstructions remove
23,078, leaving 49 associated graphs. All 49 Hamiltonian searches are
infeasible. Small separators give additional elementary explanations for
45 of them: deleting a few vertices leaves too many disconnected pieces
for a single path, sometimes after accounting for its forced endpoints.
The four remaining cases retain their native solver evidence.

This excludes every candidate whose 97-vertex majority class contains
one entire saved 96-vertex class. It allows a completely new order,
new endpoints and unrestricted opposite-parity vertices. Keeping the
96-vertex minority class is a different case: its opposite-parity
endpoint vertices need only one neighbour in the fixed class. That case
is addressed below~\cite{research231}.

\section{Why both colour classes must change}
Now keep a saved 96-vertex class as the minority class. An internal
opposite-parity vertex still needs two neighbours there, while an
endpoint needs exactly one. The contracted Hamiltonian path must start
and end at vertices admitting those endpoint labels. Inserting the
internal labels and endpoints gives $96+95+2=193$ vertices and 192 edges,
with no additional adjacency.

All 145 classes were covered. Two lack two eligible ends, 62 have an
ineligible leaf, and the other 81 queries are infeasible. In 33 of those
cases a small separator splits the augmented cycle graph into too many
components. The other 48 retain native solver evidence. Lean checked
the endpoint attachment, minority construction and supporting neighbour
statements; the finite survey remains separate.

Together with the majority result, this means every 192-edge snake must
omit at least one vertex of each parity relative to every saved seed.
Keeping a whole 96-set would put it in one of the excluded cases. The
conclusion also holds for cube-isometric copies by transforming a
proposed candidate back to the saved coordinates. It retains the native
finite-evidence qualification and gives no new global bound~\cite{research232}.
There is an exact restricted optimum as well: requiring one complete
selected colour class to equal a saved 96-set permits at most 191 edges. Alternation
allows at most 192; that case is excluded, and the saved seed attains 191.
Three further Lean counting statements check this length deduction under
its explicit count and exclusion premises.
The one-replacement route is addressed below. The intervening round
first examined a smaller family built from old paths.

\section{Why the simple dimensional lift fails}
Project the even vertices of $Q_9$ onto their first eight coordinates.
The missing bit is forced by parity, so every chosen even class is
equivalent to a subset $S\subseteq Q_8$. An odd vertex over $y$ sees
exactly the selected vertices in the closed neighbourhood $N[y]$.
Whenever $N[y]\cap S$ has exactly two vertices, join them by an edge
labelled $y$. A Hamiltonian path, which visits every vertex of $S$
once, then lifts to a snake of $2|S|-2$ edges. Extra adjacencies are
excluded by the exact two-neighbour condition. Conversely, every
even-length snake has this representation, after a coordinate flip
if needed. Thus 192 edges are equivalent to finding a suitable
97-vertex subset of $Q_8$.

One proposed source of $S$ was an old $m$-edge snake with $r$ vertices
deleted. It would give $2m-2r$ edges. With $r=1$ the arithmetic reproduces
$26\to50$ and $50\to98$, then predicts $98\to194$ if the required path
exists. The full survey of zero, one and two deletions found that every
case from the saved 26-, 50- and 98-edge parents has an isolated vertex
in its contraction. Such a vertex cannot belong to a Hamiltonian path
with other vertices. The numerical pattern therefore does not describe
this construction for those parents.

There is also a short explanation that covers more deletions. An
initially isolated vertex can change only if a vertex within Hamming
distance two is deleted; otherwise every potentially useful label sees
the same selected neighbours. For the saved 98-edge parent, 69 vertices
are initially isolated, and one deletion can affect at most 17 of them.
Four deletions cover at most $4\cdot17=68$. Hence at least five are
necessary, so this entire family has length at most $196-10=186$.
For the saved 50- and 26-edge parents, corresponding finite covering
arguments require at least four and three deletions respectively.
These are restrictions on specific construction families; they do not
give upper bounds for arbitrary snakes~\cite{research233}.

Six new general Lean statements support the lift and local-deletion
arguments. The whole-cube projection and counting bridges are written
proofs. The independent finite audit checked all 6763 graphs, 36 circuit
models and 28 successful shorter witnesses. Eight native infeasibility
reports remain separate from Lean. No longer $Q_9$ snake was found.
This motivated allowing additions as well as deletions in a saved parity
class, while leaving the opposite class and ordering free. Research
resumed after the round-233 pause.

\section{One replacement is still insufficient}
Replace one vertex of a saved 96-vertex minority class, leaving its
opposite class, endpoints and ordering free. The exact neighbour update
removes the old vertex and adds the new one wherever adjacent. A label
that sees both has the same degree but different neighbours; those
identities must be updated. Four general Lean statements check the
update and its relevant special cases.

All $2\,227\,200$ descriptions were examined and independently
reconstructed. Elementary isolated-vertex, leaf, endpoint and connectivity
conditions leave 1381 descriptions, representing 1325 distinct classes.
Of these, 1244 have independently replayed elementary contradictions.
For example, a vertex with only two available edges forces both into
the proposed spanning cycle; successive forced choices may close a
smaller cycle and prevent a single cycle through all vertices. The
remaining 81 classes retain native infeasibility evidence. Every case
is covered, with no 192-edge construction~\cite{research234}.

This sharpens the distance from every saved seed. A candidate's
96-vertex minority class can share at most 94 vertices with the
corresponding seed class: sharing 95 would be one replacement, and
sharing 96 would be the already excluded unchanged class. Its
97-vertex majority class can share at most 95 by the earlier cover.
Thus the candidate shares at most $94+95=189$ vertices with the seed.
Since the seed has 192 vertices and the candidate has 193, at least
three old vertices must be omitted and four new vertices introduced.
The argument also applies to cube-isometric seed copies, with the same
finite-evidence qualification. These changes are necessary and need
not be sufficient.

The next proposed experiment allows two replacements in a minority
class and chooses all other vertices and their ordering together.
The following round implements this approach. The global interval
remains unchanged.

\section{Two replacements also fail the degree conditions}
The next model chooses 96 minority and 97 majority vertices together,
retaining exactly 94 vertices of a saved minority class. Four direct
circuit queries remained unresolved after a minute each.

A smaller model asks only for the necessary neighbour counts. Each
selected vertex must have two selected cube neighbours, except the two
majority endpoints, which have one. Unselected vertices are unrestricted.
Any actual snake satisfies these conditions. A feasible disconnected
selection would consist of one path and some cycles; verified cycle
inequalities could then remove those cycles.

All 145 saved classes were infeasible with two replacements even in
this smaller model, before any cycle inequality was needed. The models
jointly cover $8\,410\,464\,000$ replacement descriptions. Independent
inspection verified every model and its positive selection guards.
The infeasibility conclusions remain native solver evidence. Four new
Lean statements support the separate circuit-to-path argument; the
finite contradictions are not Lean certificate replays~\cite{research235}.

The important consequence is that the candidate's minority class can
now share at most 93 vertices with the corresponding saved class.
Its majority class can share at most 95 by the earlier result, so the
total overlap is at most $93+95=188$. A saved seed has 192 vertices and
a proposed 192-edge snake has 193. At least four old vertices must
therefore be omitted and five new ones introduced. This applies to
each saved seed and its cube-isometric copies under the same finite
evidence, and does not establish a global upper bound.

The following round applies the saved degree formulation to three
replacements and investigates the conditions on larger changes.

\section{Three replacements and the next construction}
All 145 saved three-replacement families also fail the necessary degree
conditions. The exported models and the complete coverage were checked
independently. Their infeasibility remains native solver evidence.
Combining the covers through three replacements bounds the minority
intersection by 92. The majority intersection remains at most 95, so
the total overlap with any saved seed or its cube-isometric copy is at
most $92+95=187$. Thus a 192-edge candidate must omit at least five old
vertices and introduce at least six new ones.

The first four searches with exactly four replacements gave:
class 139: \texttt{INFEASIBLE}; class 7: \texttt{UNKNOWN}; class 6: \texttt{UNKNOWN}; class 0: \texttt{UNKNOWN}. These are only four saved classes; an unresolved result
does not establish impossibility.

There is a useful rule for choosing the changes. At an old internal
vertex that stays internal, every added neighbour must be balanced by
a removed neighbour. An old endpoint that becomes internal needs one
extra neighbour. This follows by subtracting its old and new degree
equations. After proposing majority-side changes, every old internal minority
vertex with unbalanced neighbour counts must be deleted, since all
retained minority vertices are internal in the candidate. This can identify required deletions before a complete search.

Nine-dimensional neighbour counts also have an exact inverse. If $b(v)$
is the change in the number of selected neighbours, the weighted sum of
all these changes around a vertex, with weights $35,-10,8,-16,128$ at
distances $1,3,5,7,9$, equals 315 times that vertex's selection change.
The proof counts how each cube step increases or decreases distance;
all coefficients cancel except the coefficient at the vertex itself.
Therefore complete neighbour counts determine the selected set uniquely.
Counts outside the retained snake are essential unknowns in this formula.

Six supporting statements passed Lean, while the full cube matrix was
checked separately by exact integer arithmetic. The geometric connection
and native exclusions have the stated separate scopes~\cite{research236}.
No longer snake or new global bound was obtained. The open construction
problem now starts with at least four minority replacements. The following round examines a complementary change on the majority
side; unresolved four-replacement searches remain saved.

\section{Changing the majority class and recovering missing counts}
The complementary construction removes one vertex from a saved
96-vertex parity set and inserts two others of the same parity. These
97 vertices would be the majority class of a 192-edge snake. The other
96 vertices and the entire ordering remain free.

Join two majority vertices when some opposite-parity vertex has exactly
those two selected neighbours. A snake gives a path visiting every
majority vertex once in this simple graph. Conversely, such an ordering
can be interleaved with its witnessing vertices to give the snake.

This gives two useful necessary conditions. If $r$ majority vertices are
still missing from a partial proposal, its isolated vertices and leaves
satisfy $2I+L\le2r+2$. Removing $r$ positions from a spanning path leaves
at most $r+1$ runs, each contributing two missing incidences; additional
graph edges can only reduce the deficit. For a complete proposal and
arbitrary subsets $S,T$, counting the two directions along the path gives
\[
2|S|\le2|T|+2+e(S,U\setminus T).
\]
Here $e$ counts ordered edges from $S$ to vertices outside $T$. Only two
directions can end at path endpoints; vertices of $T$ receive at most two.

Two independent enumerations agree on all surviving proposals among
177,062,400 descriptions. There are 679 survivors representing 623
distinct sets. Each violates the displayed cut inequality for saved
explicit subsets. All 623 numerical inequalities passed kernel checking.
The full connection from numerical counts to the path theorem remains
a written argument, with its proof-assistant assembly incomplete. A
separate audited degree model also reports all 145 starting families
infeasible. Those native exclusions have no Lean contradiction replay.

Under these finite covers and the earlier ones, the majority overlap
with a saved seed is at most 94. Combined with the minority overlap
bound 92, the total is at most $94+92=186$.
A 192-edge candidate must therefore omit at least six old vertices and
introduce at least seven new ones, relative to every saved seed or its
cube-isometric copy. This remains a restriction on proximity to the
saved collection.

The cube reconstruction theorem has also been strengthened: neighbour
counts at any 510 of the 512 vertices determine a binary selection
uniquely. To see why two missing counts cannot conceal a change, apply
the inverse formula at their antipodes. Two integer equations result,
with diagonal coefficient 128 and a short list of possible off-diagonal
coefficients. The selection difference is $-1$, zero or one. The only
integer solution consistent with these equations has both missing
count changes zero. The inverse then makes the entire selection change
zero. This actual-cube statement passed Lean after the large finite
calculation was divided into smaller modules~\cite{research237}.

No longer snake was obtained. The current interval is $191\le a(9)\le231$
using the separately cited contributor theorem. The following sections develop the reconstruction argument further.


\section{Three missing observations can also be recovered}
The recovery result now tolerates three missing counts: any 509 of the
512 neighbour counts determine the selected set. When one parity side
is known in advance, 253 of the 256 opposite-side counts suffice.

For the reasoning, compare two selections with the same observed
counts. Their selection difference is always $-1$, zero or one, and
their count difference is confined to three positions. Apply the inverse
formula at the three antipodes. This gives three integer equations with
diagonal coefficient 128 and off-diagonal coefficients among
$0,35,-10,8,-16$. The right sides are $-315$, zero or 315.
Parity and divisibility rule out every nonzero solution: the proof
separates the cases according to how many off-diagonal coefficients are
the odd number 35. Thus all three missing count differences vanish,
and the inverse makes the selections equal.

The 509-count statement and its connection to the actual cube passed
Lean. The 253-count parity specialization is a written corollary.
The main paper includes the short mathematical argument; the archive
also contains the exact coefficient calculations~\cite{research238}.
It supplies a necessary consistency condition for proposed constructions.
It does not decide whether a selection is one connected snake.
Two distinct selections must differ in at least four neighbour counts.
The following section extends this conclusion to four missing counts.
No longer snake or new global bound was obtained in this round.

\section{Why four missing counts can also be recovered}
The neighbour counts at any 508 of the 512 vertices of $Q_9$ determine
a binary vertex selection uniquely. If the selection lies on a known
parity side, any 252 of the 256 opposite-side counts suffice. The counts
on the selected side are automatically zero.

To compare two selections, subtract their indicators. Their difference
$f$ takes values $-1,0,1$. If the observed counts agree, the neighbour
sums of $f$ can be nonzero only at the four missing positions. When
those positions have mixed parity, splitting $f$ by parity reduces
each part to at most three missing counts, already covered by the
preceding result.

For four positions of one parity, apply the inverse formula at their
antipodes. If $w$ lists the four missing count differences, the resulting
integer equations are
\[
Mw=315s,\qquad s\in\{-1,0,1\}^4.
\]
The diagonal entries of $M$ are 128. Its other entries are $-16,8,-10,35$
when the corresponding missing positions are at distance $2,4,6,8$.

Only 35 is odd. Join positions at distance eight. Two such neighbours
of one position have distance two from each other, so this graph has
no triangle. There are seven possible graph types on four vertices.
Parity, subtraction of equations and divisibility settle five types
directly. The other two reduce to seven triples of coefficients and
five small block matrices. In each case, elimination produces
$Dx=315N$; the integer $D/\gcd(D,315)$ must divide $N$, while the allowed
right-hand sides make $|N|$ too small unless $N=0$. The full paper gives
the matrices and numerical bounds. These arguments force $w=0$, and
the inverse then forces the entire selection difference $f$ to vanish.

A further result permits up to eight missing counts when no pair of
missing positions has Hamming distance eight. In that case every
off-diagonal coefficient in the antipodal equations is even, with
absolute value at most 16. Parity forces every right-hand side to zero.
If $k\le8$ counts are missing and $M$ is the largest absolute missing
count difference, a row attaining it gives
\[
128M\le16(k-1)M\le112M,
\]
so $M=0$. The supporting integer inequalities passed Lean; their
application to the cube remains a written deduction. The condition
on the missing positions is required by this argument.

Eight is sharp under that condition. The nine neighbours of one vertex
have pairwise distance two, yet losing all nine counts conceals whether
that vertex is selected. The unrestricted result above remains four
missing counts; whether every set of five through eight missing counts
can be tolerated is unresolved. Any failure in that range must contain
a pair of missing positions at distance eight.

The classification was also checked independently. Grouping the nine
coordinate columns by their three-bit patterns gives 1149 valid
ordered descriptions, 105 cube-isometry classes of four-point sets,
and 87 distinct distance matrices. A second enumeration checks all
2,731,135 triples after fixing one vertex at zero and obtains exactly
the same 87 matrices. All 7047 matrix/right-hand-side combinations have
no nonzero integer solution~\cite{research239}.

All 87 matrix implications passed Lean checking, including their
adjugate identities and divisibility calculations. The connection from
arbitrary cube positions to this catalogue, and the parity reduction,
remain written arguments with independent computational support. The
complete four-count recovery statement has not been assembled as one
Lean theorem.

Thus two different selections must alter at least five neighbour
counts. Changing one selected vertex alters nine, so the smallest
possible number remains between five and nine. Recovery alone does
not ensure that a selection is a connected induced path. No longer
snake or new global bound follows from this result.

\section{Remaining questions}
A one-crossing join using the full saved parent class cannot improve 191.
The selected two-block family is excluded through width 46 for every
crossing count. A new construction search must change a restriction.
The parity-class comparisons exclude keeping either entire class of a
saved seed. Changes to both classes, wider blocks and changed block
lengths remain possible.
Different parent classes and shorter compatible parents remain alternatives. Finding a new parent still requires a compatible joining construction.
Exhausting all remaining six-coordinate models would establish only a
source-specific requirement of seven changed histories.
The maintained global interval is $191\le a(9)\le231$, with an
independent rebuild of the completed upper-bound assembly pending.
Resolving the proposed value 196 still requires both its witness and
a universal 197-edge exclusion. For reconstruction, the next unresolved
case has five missing positions of one parity containing a distance-eight
pair. Fixing that pair by cube symmetry reduces a complete description
cover to $\binom{15}{7}=6435$ coordinate-column multiplicities. This
reduction is proposed for the next investigation; it has not been searched.
Research is paused after round 239.

\smallskip
{\small Detailed proofs, witnesses and verification records are retained
in the companion manuscript and milestones through round 239.
The small-dimensional values appear in
\href{https://oeis.org/A099155}{OEIS A099155}; the published 191-edge
construction is in
\href{https://arxiv.org/html/2607.15270v3}{Orland et al., arXiv:2607.15270v3}.}

\begin{thebibliography}{9}
\bibitem{bound231archive}
Snake-in-the-box project contributors. \emph{Completed upper bound 231:
final assembly and standard-axiom audit}. Contributor revision
\texttt{478869ee1ec6ff3a10ddb6ba7a2a1d89fbb82088}, 18 November 2026.
A fresh independent whole-proof rebuild remains pending.
\url{https://github.com/hironaka-iori/Snake-in-the-box-in-9/tree/478869ee1ec6ff3a10ddb6ba7a2a1d89fbb82088/milestones/36-lean435-port/audit231}.

\bibitem{bound233archive}
Snake-in-the-box project contributors. \emph{Completed share-weighted bound
233: final assembly, standard-axiom audit and verification limitations}.
Contributor revision \texttt{b0521e1d14fd52c62db4cdeccc356b70a623eb82};
audit revision \texttt{b65d30ca15441420b8f4538f671e589c2502674f},
14 November 2026.
\url{https://github.com/hironaka-iori/Snake-in-the-box-in-9/tree/b0521e1d14fd52c62db4cdeccc356b70a623eb82/milestones/33-233-lean}.


\bibitem{research220}
I.~Hironaka, T.~Ma, and S.~Zhan. \emph{Subset forest bounds,
fractional feasibility and product consistency}. Research milestone 220,
15 November 2026.
\url{https://github.com/hironaka-iori/Snake-in-the-box-in-9/tree/main/milestones/220-subset-forest}.

\bibitem{research221}
I.~Hironaka, T.~Ma, and S.~Zhan. \emph{Cycle parity, exact partner bounds
and first-step restrictions}. Research milestone 221, 15 November 2026.
\url{https://github.com/hironaka-iori/Snake-in-the-box-in-9/tree/main/milestones/221-square-parity}.

\bibitem{research222}
I.~Hironaka, T.~Ma, and S.~Zhan. \emph{Complete prefix certificates
and a one-crossing ceiling for the saved parent class}.
Research milestone 222, 15 November 2026.
\url{https://github.com/hironaka-iori/Snake-in-the-box-in-9/tree/main/milestones/222-prefix-certificates}.

\bibitem{barahona1986}
F.~Barahona and A.~R.~Mahjoub. On the cut polytope.
\emph{Mathematical Programming} \textbf{36} (1986), 157--173.
\url{https://www.lamsade.dauphine.fr/~mahjoub/Papers/BM86.pdf}.

\bibitem{research223}
I.~Hironaka, T.~Ma, and S.~Zhan. \emph{Joint four-crossing construction
and five-component joining}. Research milestone 223, 15 November 2026.
\url{https://github.com/hironaka-iori/Snake-in-the-box-in-9/tree/main/milestones/223-coupled-four-crossings}.

\bibitem{research224}
I.~Hironaka, T.~Ma, and S.~Zhan. \emph{Interaction between simultaneous
replacements and two-region construction searches}.
Research milestone 224, 15 November 2026.
\url{https://github.com/hironaka-iori/Snake-in-the-box-in-9/tree/main/milestones/224-cross-block-reconnections}.

\bibitem{research225}
I.~Hironaka, T.~Ma, and S.~Zhan. \emph{Simultaneous replacement criteria,
containment transfer and crossing-parity search}. Research milestone 225,
15 November 2026.
\url{https://github.com/hironaka-iori/Snake-in-the-box-in-9/tree/main/milestones/225-replacement-widths}.

\bibitem{research226}
I.~Hironaka, T.~Ma, and S.~Zhan. \emph{Residual capacity, boundary changes
and shared eight-dimensional replacement paths}. Research milestone 226,
15 November 2026.
\url{https://github.com/hironaka-iori/Snake-in-the-box-in-9/tree/main/milestones/226-connected-replacement-search}.

\bibitem{research227}
I.~Hironaka, T.~Ma, and S.~Zhan. \emph{Fixed upper portions and exact
lower-dimensional replacement tasks}. Research milestone 227,
15 November 2026.
\url{https://github.com/hironaka-iori/Snake-in-the-box-in-9/tree/main/milestones/227-shared-q8-construction}.

\bibitem{research228}
I.~Hironaka, T.~Ma, and S.~Zhan. \emph{Two lower-dimensional replacement
paths and the width-45 construction cover}. Research milestone 228,
15 November 2026.
\url{https://github.com/hironaka-iori/Snake-in-the-box-in-9/tree/main/milestones/228-wider-replacement-constructions}.

\bibitem{research229}
I.~Hironaka, T.~Ma, and S.~Zhan. \emph{Exact projection restrictions and
variable-crossing replacement constructions}. Research milestone 229,
15 November 2026.
\url{https://github.com/hironaka-iori/Snake-in-the-box-in-9/tree/main/milestones/229-two-layer-construction}.

\bibitem{research230}
I.~Hironaka, T.~Ma, and S.~Zhan. \emph{Face visits and completion of the
two-block construction exclusion}. Research milestone 230,
16 November 2026.
\url{https://github.com/hironaka-iori/Snake-in-the-box-in-9/tree/main/milestones/230-two-crossing-frontier}.

\bibitem{research231}
I.~Hironaka, T.~Ma, and S.~Zhan. \emph{Parity-class contraction and
Hamiltonian construction searches}. Research milestone 231,
16 November 2026.
\url{https://github.com/hironaka-iori/Snake-in-the-box-in-9/tree/main/milestones/231-endpoint-exchanges}.

\bibitem{research232}
I.~Hironaka, T.~Ma, and S.~Zhan. \emph{Minority-class contraction and
the need to change both parity classes}. Research milestone 232,
16 November 2026.
\url{https://github.com/hironaka-iori/Snake-in-the-box-in-9/tree/main/milestones/232-minority-contraction}.

\bibitem{research233}
I.~Hironaka, T.~Ma, and S.~Zhan. \emph{Parity lift and deletion obstruction}.
Research milestone 233, 16 November 2026.
\url{https://github.com/hironaka-iori/Snake-in-the-box-in-9/tree/main/milestones/233-parity-lift-construction}.

\bibitem{research234}
I.~Hironaka, T.~Ma, and S.~Zhan. \emph{Replacing one member of a parity class}.
Research milestone 234, 17 November 2026.
\url{https://github.com/hironaka-iori/Snake-in-the-box-in-9/tree/main/milestones/234-single-colour-exchange}.

\bibitem{research235}
I.~Hironaka, T.~Ma, and S.~Zhan. \emph{Joint construction and the
two-replacement obstruction}. Research milestone 235, 17 November 2026.
\url{https://github.com/hironaka-iori/Snake-in-the-box-in-9/tree/main/milestones/235-joint-two-exchange}.

\bibitem{research236}
I.~Hironaka, T.~Ma, and S.~Zhan. \emph{Larger parity exchanges and
neighbour-count reconstruction}. Research milestone 236, 18 November 2026.
\url{https://github.com/hironaka-iori/Snake-in-the-box-in-9/tree/main/milestones/236-parity-exchange-frontier}.

\bibitem{research237}
I.~Hironaka, T.~Ma, and S.~Zhan. \emph{Majority exchanges and reconstruction
with missing neighbour counts}. Research milestone 237, 18 November 2026.
\url{https://github.com/hironaka-iori/Snake-in-the-box-in-9/tree/main/milestones/237-parity-reconstruction}.

\bibitem{research238}
I.~Hironaka, T.~Ma, and S.~Zhan. \emph{Recovery from three missing neighbour
counts}. Research milestone 238, 18 November 2026.
\url{https://github.com/hironaka-iori/Snake-in-the-box-in-9/tree/main/milestones/238-three-count-recovery}.

\bibitem{research239}
I.~Hironaka, T.~Ma, and S.~Zhan. \emph{Four missing neighbour counts and a
conditional eight-count theorem}. Research milestone 239, 18 November 2026.
\url{https://github.com/hironaka-iori/Snake-in-the-box-in-9/tree/main/milestones/239-four-count-recovery}.

\end{thebibliography}
\end{document}
